% GENERATED FILE. Edit canonical lessons, then run npm run generate:books.
% canonical-source-hash: fnv1a64:896df357b073d273
% canonical-lessons: KA-C245-bhedi, KA-C245-alarji, KA-C245-uta, KA-C245-gulle, KA-C245-cikitse

\chapter{Diarrhoea, Allergy, Swelling, Pimple, Medical treatment}
\label{ch:ka-diarrhoea-allergy-swelling-pimple-medical-treatment}
\hlchaptermodality{245}

\begin{chapteropening}
\textbf{By the end of this chapter:} \emph{I can say diarrhoea, an allergy, a swelling, a pimple and medical treatment.}

Complete the last lesson of chapter 245: I can say diarrhoea, an allergy, a swelling, a pimple and medical treatment.
\end{chapteropening}

\section[\texorpdfstring{\kn{ಭೇದಿ} (bhēdi)}{ಭೇದಿ (bhēdi)}]{\kn{ಭೇದಿ} (bhēdi) — diarrhoea}

\label{lesson:KA-C245-bhedi}



\begin{quote}
Before the new one: say the Kannada for to suck, then the Kannada for to install.
\end{quote}

\subsection*{You'll want to know: \kn{ಭೇದಿ}}
\textbf{\kn{ಭೇದಿ}} — \emph{bhēdi} — \textquotedblleft{}diarrhoea\textquotedblright{}.

\begin{grammarlens}[title={What you've built}]
One more word for this chapter.
\end{grammarlens}

\subsection*{Your turn}
\begin{itemize}
\raggedright
  \item \emph{Say it:} \emph{bhēdi}
  \item \emph{Say it:} \emph{bhēdi}, once more
  \item \emph{Say it:} say \emph{bhēdi}
  \item \emph{Recall:} say the Kannada for to suck, then the Kannada for to install, then say \emph{bhēdi} again
\end{itemize}

\begin{tcolorbox}[breakable,colback=teal!4,colframe=teal!35!black,title={Before you move on}]
What does \kn{ಭೇದಿ} mean? (\textquotedblleft{}Diarrhoea\textquotedblright{}.) Say it once more.
\end{tcolorbox}
\section[\texorpdfstring{\kn{ಅಲರ್ಜಿ} (alarji)}{ಅಲರ್ಜಿ (alarji)}]{\kn{ಅಲರ್ಜಿ} (alarji) — an allergy}

\label{lesson:KA-C245-alarji}



\begin{quote}
Before the new one: say the Kannada for to install, then the Kannada for diarrhoea.
\end{quote}

\subsection*{You'll want to know: \kn{ಅಲರ್ಜಿ}}
\textbf{\kn{ಅಲರ್ಜಿ}} — \emph{alarji} — \textquotedblleft{}an allergy\textquotedblright{}.

\begin{grammarlens}[title={What you've built}]
One more word for this chapter.
\end{grammarlens}

\subsection*{Your turn}
\begin{itemize}
\raggedright
  \item \emph{Say it:} \emph{alarji}
  \item \emph{Say it:} \emph{alarji}, once more
  \item \emph{Say it:} say \emph{alarji}
  \item \emph{Recall:} say the Kannada for to install, then the Kannada for diarrhoea, then say \emph{alarji} again
\end{itemize}

\begin{tcolorbox}[breakable,colback=teal!4,colframe=teal!35!black,title={Before you move on}]
What does \kn{ಅಲರ್ಜಿ} mean? (\textquotedblleft{}An allergy\textquotedblright{}.) Say it once more.
\end{tcolorbox}
\section[\texorpdfstring{\kn{ಊತ} (ūta)}{ಊತ (ūta)}]{\kn{ಊತ} (ūta) — a swelling}

\label{lesson:KA-C245-uta}



\begin{quote}
Before the new one: say the Kannada for diarrhoea, then the Kannada for an allergy.
\end{quote}

\subsection*{You'll want to know: \kn{ಊತ}}
\textbf{\kn{ಊತ}} — \emph{ūta} — \textquotedblleft{}a swelling\textquotedblright{}.

\begin{grammarlens}[title={What you've built}]
One more word for this chapter.
\end{grammarlens}

\subsection*{Your turn}
\begin{itemize}
\raggedright
  \item \emph{Say it:} \emph{ūta}
  \item \emph{Say it:} \emph{ūta}, once more
  \item \emph{Say it:} say \emph{ūta}
  \item \emph{Recall:} say the Kannada for diarrhoea, then the Kannada for an allergy, then say \emph{ūta} again
\end{itemize}

\begin{tcolorbox}[breakable,colback=teal!4,colframe=teal!35!black,title={Before you move on}]
What does \kn{ಊತ} mean? (\textquotedblleft{}A swelling\textquotedblright{}.) Say it once more.
\end{tcolorbox}
\section[\texorpdfstring{\kn{ಗುಳ್ಳೆ} (guḷḷe)}{ಗುಳ್ಳೆ (guḷḷe)}]{\kn{ಗುಳ್ಳೆ} (guḷḷe) — a pimple, a blister}

\label{lesson:KA-C245-gulle}



\begin{quote}
Before the new one: say the Kannada for an allergy, then the Kannada for a swelling.
\end{quote}

\subsection*{You'll want to know: \kn{ಗುಳ್ಳೆ}}
\textbf{\kn{ಗುಳ್ಳೆ}} — \emph{guḷḷe} — \textquotedblleft{}a pimple, a blister\textquotedblright{}.

\begin{grammarlens}[title={What you've built}]
One more word for this chapter.
\end{grammarlens}

\subsection*{Your turn}
\begin{itemize}
\raggedright
  \item \emph{Say it:} \emph{guḷḷe}
  \item \emph{Say it:} \emph{guḷḷe}, once more
  \item \emph{Say it:} say \emph{guḷḷe}
  \item \emph{Recall:} say the Kannada for an allergy, then the Kannada for a swelling, then say \emph{guḷḷe} again
\end{itemize}

\begin{tcolorbox}[breakable,colback=teal!4,colframe=teal!35!black,title={Before you move on}]
What does \kn{ಗುಳ್ಳೆ} mean? (\textquotedblleft{}A pimple, a blister\textquotedblright{}.) Say it once more.
\end{tcolorbox}
\section[\texorpdfstring{\kn{ಚಿಕಿತ್ಸೆ} (cikitse)}{ಚಿಕಿತ್ಸೆ (cikitse)}]{\kn{ಚಿಕಿತ್ಸೆ} (cikitse) — medical treatment}

\label{lesson:KA-C245-cikitse}



\begin{quote}
Before the new one: say the Kannada for a swelling, then the Kannada for a pimple.
\end{quote}

\subsection*{You'll want to know: \kn{ಚಿಕಿತ್ಸೆ}}
\textbf{\kn{ಚಿಕಿತ್ಸೆ}} — \emph{cikitse} — \textquotedblleft{}medical treatment\textquotedblright{}.

\begin{grammarlens}[title={What you've built}]
One more word for this chapter.
\end{grammarlens}

\subsection*{Your turn}
\begin{itemize}
\raggedright
  \item \emph{Say it:} \emph{cikitse}
  \item \emph{Say it:} \emph{cikitse}, once more
  \item \emph{Say it:} say \emph{cikitse}
  \item \emph{Recall:} say the Kannada for a swelling, then the Kannada for a pimple, then say \emph{cikitse} again
\end{itemize}

\begin{tcolorbox}[breakable,colback=teal!4,colframe=teal!35!black,title={Before you move on}]
What does \kn{ಚಿಕಿತ್ಸೆ} mean? (\textquotedblleft{}Medical treatment\textquotedblright{}.) Say it once more.
\end{tcolorbox}
